\documentclass[11pt]{article}
\usepackage{fontspec}
\setmainfont{Latin Modern Roman}
\usepackage{amsmath,amssymb,mathtools,unicode-math}
\setmathfont{Latin Modern Math}
\usepackage{xurl,hyperref}
\hypersetup{hidelinks,pdftitle={Problem statement}}
\newcommand{\sref}[2]{\hyperlink{source:#1:#2}{[S#2]}}
\newcommand{\cataloguescope}[1]{\paragraph{Mathematical question.}#1}
\begin{document}
\section{Open Problems in Integrable Systems — Two-dimensional case.}\label{statement}

\subsection{Definitions and mathematical statement}

\paragraph{Shared notation.}$\mathbb N=\{1,2,\ldots\}$, and $\mathbb Z,\mathbb Q,\mathbb R,\mathbb C$ denote the integers, rationals, reals and complex numbers. Any other conventions are specified locally.


For a smooth Riemannian metric $g$ on $M$, a natural Hamiltonian on the canonical symplectic cotangent bundle $T^*M$ is \[H(x,p)=\tfrac12\sum_{i,j}g^{ij}(x)p_ip_j+U(x),\] with smooth potential $U$ and inverse metric coefficients $g^{ij}$. A polynomial integral of degree $d$ is a smooth function $F(x,p)$ polynomial of degree $d$ in the cotangent variables $p$, satisfying $\{H,F\}=0$. Following the source, integrals which are functions of $H$ alone are trivial. A geodesic flow is the case $U=0$.
\paragraph{Statement recorded in the supplied source.}
Describe and classify the natural systems on closed two-dimensional surfaces admitting a nontrivial polynomial integral, organized by the smallest such degree: construct examples and prove nonexistence results, with any additional hypotheses made explicit. The source's table covers $S^2,T^2$ and degrees $1,2,3,4,\ge5$. Its precise subsidiary conjectures ask: (i) whether a cubic polynomial integral for a Riemannian geodesic flow on $T^2$ forces a Killing vector field; (ii) whether, on $T^2$ with nonconstant $U$, an odd-degree polynomial integral forces a linear integral and a nontrivial even-degree integral forces a nontrivial quadratic integral; and (iii) the same assertion in (ii) when the torus metric is flat. A Killing vector field generates local isometries and yields a linear integral. \sref{statement}{2}
\subsection{Short English statement}
Classify polynomially integrable natural systems on closed surfaces by their smallest integral degree, including the stated torus reductions to linear or quadratic integrals.
\subsection{Sources}
\begin{description}
\item[\hypertarget{source:statement:1}{S1}] ULAM.AI and the UnsolvedMath Contributors, UnsolvedMath: A Curated Collection of Open Mathematics Problems, catalogue entry. CC BY 4.0. \url{https://huggingface.co/datasets/ulamai/UnsolvedMath}.
\item[\hypertarget{source:statement:2}{S2}] A. Bolsinov, V. S. Matveev, E. Miranda and S. Tabachnikov, Open Problems, Questions, and Challenges in Finite-Dimensional Integrable Systems (2018), arXiv:1804.03737, Problem 3.1, table, and Conjectures 3.2–3.4. \url{https://arxiv.org/abs/1804.03737}.
\end{description}
\label{end:statement}
\end{document}
